\section{Sobre o Autor}

\begin{figure}[h]
  \centering

  \begin{minipage}{0.28\textwidth}
    \centering
    \includegraphics[width=\linewidth]{./fig/Augusto-Frederico-Schmidt.jpg}

    \vspace{0.5em}
    \footnotesize
    Augusto Frederico Schmidt\\
    (1906--1965)
  \end{minipage}
  \hfill
  \begin{minipage}{0.7\textwidth}
    Augusto Frederico Schmidt (Rio de Janeiro, 18 de abril de 1906 -- Rio
    de Janeiro, 8 de fevereiro de 1965) foi poeta, editor, empresário,
    diplomata e homem público brasileiro. Filho de um imigrante judeu
    asquenaze alemão e neto do Visconde de Schmidt, iniciou sua trajetória
    no mercado editorial após atuar na Livraria Católica e, em 1930, fundou
    a Livraria Schmidt Editora.\\

    À frente da editora, publicou as primeiras obras de autores que se
    tornariam fundamentais para a literatura brasileira, entre eles Gilberto
    Freyre, Graciliano Ramos, Rachel de Queiroz, Jorge Amado, Vinicius de
    Moraes, Lúcio Cardoso e Octávio de Faria. A Livraria Schmidt tornou-se um
    dos principais centros da vida intelectual carioca na década de 1930,
    servindo como ponto de encontro do chamado \textit{Círculo Católico},
    grupo informal que reunia escritores, jornalistas e pensadores ligados ao
    renascimento do pensamento católico no Brasil na época. \\

    Paralelamente à atividade literária, desenvolveu uma carreira de
    sucesso como empresário, atuando nos setores editorial, de seguros,
    mineração, aviação e varejo. Também exerceu funções diplomáticas e foi
    assessor especial do presidente Juscelino Kubitschek, sendo
    tradicionalmente associado à criação do slogan ``50 anos em 5''.
    Faleceu no Rio de Janeiro em 1965, deixando uma obra poética que
    permanece entre as mais singulares da literatura brasileira do século
    XX.
  \end{minipage}
\end{figure}

\section{Por que este poema?}

A escolha de \textit{O Pássaro} dá continuidade à análise de \textit{Era um
grande pássaro}, poema de Augusto Frederico Schmidt já estudado anteriormente
\cite{schmitz2026eraumgrandepassaro}. Nos dois textos, a figura do pássaro
ocupa uma posição central, mas assume sentidos bastante distintos.

Em \textit{Era um grande pássaro}, a imagem da ave está associada à morte, ao
mistério e à cruz. Em \textit{O Pássaro}, por outro lado, o pássaro aparece
como figura de movimento, nascimento e transformação. Essa diferença torna os
dois poemas especialmente interessantes quando lidos em conjunto, pois permite
observar como Schmidt utiliza uma mesma imagem para construir experiências
poéticas distintas.

A escolha de \textit{O Pássaro} também se justifica pela forte presença de
imagens ligadas à criação: trevas, águas, luz, vida e tempo. Essas imagens
estabelecem uma relação particularmente significativa com o relato da Criação
em Gênesis, que será o principal eixo da análise a seguir.
